\section{Three Fields, Three Failures}
\label{sec:ch15-three-fields-three-failures}

\index{resolvability!what the approved graph does at run time|(}%
The config that came out of the last section is a file a router will load, so
start the three services and point one at it. Ask for the first of the three
fields, by the first of the six routes.

\begin{minted}{graphql}
{ sessions(first: 10) { nodes { title averageScore } } }
\end{minted}

\begin{minted}[fontsize=\footnotesize,breakanywhere]{json}
{"errors":[{"message":"internal server error"}]}
\end{minted}

\index{HTTP 500!the first one in this book}%
Five hundred, and nothing in this book has answered one before. Every response
chapter~\ref{ch:null-blast-radius} printed was a two hundred with the problem
in the envelope, and the one refusal chapter~\ref{ch:entities-authorization}
met was a four hundred and one from the router about a token. This is neither.
The router took a request it advertises, planned it, and could not.

\index{HTTP 500!on every route}%
All six routes answer the same way, the mutation included. The router's log
says why, and it says it twice in two different vocabularies. For the four
routes that reach \code{Session} through an ordinary field:

\begin{minted}[fontsize=\footnotesize,breakanywhere]{text}
printOperation planner id: 1: validation failed: external: Cannot query field
"averageScore" on type "Query"., locations: [], path: [query]
\end{minted}

\index{router!the planner falls back to the root type}%
With no entity route into the Ratings service, the planner built the only other
kind of request it knows how to send to a subgraph: a query on that subgraph's
\code{Query} type. There is no \code{Query.averageScore}, so its own validation
caught the operation before it went out. For the two routes through
\code{node}, the log says something else:

\begin{minted}[fontsize=\footnotesize,breakanywhere]{text}
could not resolve a field: internal: nodesResolvableVisitor: could not select
the datasource to resolve Query.node on path query.node
\end{minted}

\index{router!has a resolvability check of its own}%
The router has a resolvability check of its own, it runs at request time, and
here it is finding exactly what the composer's did not.

\subsection{The Field That Answers the Wrong Question}

\index{ratingCount@\code{ratingCount}!a valid query against the wrong field}%
Ask for \code{ratingCount} instead. Same service, same defect, same route, and
now a two hundred:

\begin{minted}[fontsize=\footnotesize,breakanywhere]{json}
{"errors":[{"message":"Failed to fetch from Subgraph 'ratings' at Path 'sessions.nodes', Reason: no data or errors in response.","extensions":{"statusCode":200}}],
"data":{"sessions":{"nodes":[
{"title":"Schemas That Outlive Their Authors","ratingCount":null},
{"title":"Reading a Query Plan","ratingCount":null},
{"title":"Paging Without Offsets","ratingCount":null},
{"title":"The Bank That Split Its Graph","ratingCount":null}]}}}
\end{minted}

\index{query plan!the root field sent at a list position}%
The plan is where this gets interesting. The second step of it, with the
fields that are not the point removed:

\begin{minted}[fontsize=\footnotesize,breakanywhere]{json}
{"kind":"Single","path":"sessions.nodes","subgraphName":"ratings",
 "query":"{\n    ratingCount\n}"}
\end{minted}

\index{ratingCount@\code{ratingCount}!the collision is the whole difference}%
The router sent \code{\{ ratingCount \}} to the Ratings service: the root
field, the one chapter~\ref{ch:null-blast-radius} made nullable, at a path
inside a list of sessions. It is the same fallback that produced the five
hundred, and this time it validated, because the Ratings service happens to
have a root field with the name of the one the router wanted. A name collision
between a root field and a field on a type is the entire difference between a
loud failure and a query that runs and answers about the wrong thing.

\subsection{The One That Says Nothing}

\index{feedbackUrl@\code{feedbackUrl}!four nulls and no error}%
The third field is the quiet one:

\begin{minted}[fontsize=\footnotesize,breakanywhere]{json}
{"data":{"sessions":{"nodes":[
{"title":"Schemas That Outlive Their Authors","feedbackUrl":null},
{"title":"Reading a Query Plan","feedbackUrl":null},
{"title":"Paging Without Offsets","feedbackUrl":null},
{"title":"The Bank That Split Its Graph","feedbackUrl":null}]}}}
\end{minted}

\index{feedbackUrl@\code{feedbackUrl}!no errors key at all}%
Two hundred, every url null, and no \code{errors} key. Not an empty one. There
is nothing in the response to tell a client that anything happened.

\index{requires@\code{"@requires}!plans an entity fetch with no entity route}%
\code{feedbackUrl} is the field chapter~\ref{ch:second-third-subgraph} put
behind \code{@requires}, and that is why it behaves differently: the router
plans a \code{BatchEntity} step for it, carrying an \code{\_entities} query and
a fragment collecting \code{title}, and sends that to a service whose exported
schema has no \code{\_entities} field on it. The fetch comes back with nothing
usable and the position is nullable, so the position gets a null.

\index{resolvability!one defect, three failure modes}%
One defect. One line of C\#. Three fields on one type in one service, and three
different things happen: a five hundred, a valid query against a different
field, and silence. What sorts them is not the mistake, which is identical for
all three. It is whether a root field of the same name exists to fall back to,
and whether the field carries a \code{@requires}.

\index{composition!approved all of it}%
And a composition step approved all of it, at exit zero, with nothing on
standard output.
\index{resolvability!what the approved graph does at run time|)}%
